\begin{abstract}
Clean-label backdoor attacks are appealing because they leave poison-label
annotations untouched, yet current methods become unreliable when the victim
task is hard---many-class, low-poison-rate, or high-base-accuracy regimes where
a small, fixed trigger is diluted by the strong benign signal.
We present \textbf{FAAT} (Feature-Aligned Adaptive Trigger), a clean-label
attack composed of three generalized components: (i) a \emph{frozen global
direction} $\dglobal$ that supplies a strong, transferable backdoor signal;
(ii) a \emph{bounded adaptive residual} $\dadapt$ shaped in the DCT domain that
restores effectiveness while staying imperceptible; and (iii) a
\emph{feature-alignment loss} $\Lalign$ that pulls poisoned representations
toward the target class so sample-clustering defenses cannot isolate them.
FAAT is trained self-containedly, with no external pre-optimized noise.
On CIFAR-10, CIFAR-100 and Tiny-ImageNet, FAAT raises attack success rate
by $+3.9$ to $+51.9$ points over the strongest reproduced clean-label baseline,
with benign accuracy intact (SSIM $0.92$--$0.98$, $\ell_2\!\le\!1.5$ on natural
images), and is not flagged by Activation Clustering, Spectral Signatures,
STRIP or Fine-Pruning. As a complementary contribution we introduce \textbf{KST},
a trigger whose spectrum is \emph{structurally flat} (peak/median $\approx\!1.0$,
vs.\ $68.8$ for Narcissus), defeating frequency-domain defenses at ASR parity.
Finally, we treat GTSRB not as a fourth dataset but as a \emph{stress test of
overconfident victims} (BA$\approx\!100\%$): there, fixed, flat-spectrum and
adaptive triggers each reveal a distinct failure mode, exposing a fundamental
ASR--stealth trade-off that tightens as victim confidence grows.
\end{abstract}
